% GENERATED FILE. Edit canonical lessons, then run npm run generate:books.
% canonical-source-hash: fnv1a64:23d5c189346090a1
% canonical-lessons: KA-C506-hode-cover, KA-C506-kanavarisu, KA-C506-horalu, KA-C506-maimuri, KA-C506-reppe

\chapter{Sleepy Nights}
\label{ch:ka-sleepy-nights}
\hlchaptermodality{506}

\begin{chapteropening}
\textbf{By the end of this chapter:} \emph{I can say to cover oneself, to talk in sleep, to roll over, to stretch and an eyelid.}

Complete the last lesson of chapter 506: I can say to cover oneself, to talk in sleep, to roll over, to stretch and an eyelid.
\end{chapteropening}

\section[\texorpdfstring{\kn{ಹೊದೆ} (hode)}{ಹೊದೆ (hode)}]{\kn{ಹೊದೆ} (hode) — to cover oneself; to wrap}

\label{lesson:KA-C506-hode-cover}



\begin{quote}
Before the new one: say the Kannada for to strive, then the Kannada for to beg.
\end{quote}

\subsection*{You'll want to know: \kn{ಹೊದೆ}}
\textbf{\kn{ಹೊದೆ}} — \emph{hode} — \textquotedblleft{}to cover oneself; to wrap\textquotedblright{}.

\begin{quote}
\kn{ಚಳಿಗೆ} \kn{ಅಜ್ಜ} \kn{ಕಂಬಳಿ} \kn{ಹೊದ್ದು} \kn{ಮಲಗಿದರು}.
\end{quote}

\begin{quote}
\emph{caḷige ajja kambaḷi hoddu malagidaru.}
\end{quote}

\begin{quote}
In the cold, grandfather slept wrapped in a blanket.
\end{quote}

Past (I …): \kn{ಹೊದ್ದೆ}, \emph{hodde}.

The participle is \kn{ಹೊದ್ದು}: \kn{ಕಂಬಳಿ} \kn{ಹೊದ್ದು}, wrapped in a blanket. \kn{ಹೊದ್ದುಕೊ}, to wrap yourself up. Keep it apart from \kn{ಹೋದೆ}, I went, with a long ō.

\begin{cousinweb}
\textbf{\kn{ಹೊದೆ}} (\emph{hode}) is native Dravidian: DEDR 4509 lists Kannada \emph{pode}, \textquotedblleft{}to put over, put on, wrap round\textquotedblright{}, and \emph{podisu}, \textquotedblleft{}to cover\textquotedblright{}. Old p- became h-.
\end{cousinweb}

\begin{grammarlens}[title={What you've built}]
One more word for describing what you lived and felt.
\end{grammarlens}

\subsection*{Your turn}
\begin{itemize}
\raggedright
  \item \emph{Say it:} \emph{hode}
  \item \emph{Say it:} \emph{hode}, once more
  \item \emph{Say it:} \emph{caḷige ajja kambaḷi hoddu malagidaru.} — In the cold, grandfather slept wrapped in a blanket.
  \item \emph{Recall:} say the Kannada for to strive, then the Kannada for to beg, then say \emph{hode} again
\end{itemize}

\begin{tcolorbox}[breakable,colback=teal!4,colframe=teal!35!black,title={Before you move on}]
What does \kn{ಹೊದೆ} mean? (\textquotedblleft{}To cover oneself; to wrap\textquotedblright{}.) Say it once more.
\end{tcolorbox}
\section[\texorpdfstring{\kn{ಕನವರಿಸು} (kanavarisu)}{ಕನವರಿಸು (kanavarisu)}]{\kn{ಕನವರಿಸು} (kanavarisu) — to talk in sleep}

\label{lesson:KA-C506-kanavarisu}



\begin{quote}
Before the new one: say the Kannada for to beg, then the Kannada for to cover oneself.
\end{quote}

\subsection*{You'll want to know: \kn{ಕನವರಿಸು}}
\textbf{\kn{ಕನವರಿಸು}} — \emph{kanavarisu} — \textquotedblleft{}to talk in sleep\textquotedblright{}.

\begin{quote}
\kn{ನಿನ್ನೆ} \kn{ರಾತ್ರಿ} \kn{ತಮ್ಮ} \kn{ನಿದ್ದೆಯಲ್ಲಿ} \kn{ಕನವರಿಸಿದನು}.
\end{quote}

\begin{quote}
\emph{ninne rātri tamma niddeyalli kanavarisidanu.}
\end{quote}

\begin{quote}
Last night my younger brother talked in his sleep.
\end{quote}

Past (I …): \kn{ಕನವರಿಸಿದೆ}, \emph{kanavariside}.

Built on \kn{ಕನಸು}, a dream: talking from inside a dream.

\begin{cousinweb}
\textbf{\kn{ಕನವರಿಸು}} (\emph{kanavarisu}) is native Dravidian: DEDR 1407 lists Kannada \emph{kana}, \emph{kanasu}, \textquotedblleft{}a dream\textquotedblright{}, and \emph{kanavarisu}, a verb of dreaming.
\end{cousinweb}

\begin{grammarlens}[title={What you've built}]
One more word for describing what you lived and felt.
\end{grammarlens}

\subsection*{Your turn}
\begin{itemize}
\raggedright
  \item \emph{Say it:} \emph{kanavarisu}
  \item \emph{Say it:} \emph{kanavarisu}, once more
  \item \emph{Say it:} \emph{ninne rātri tamma niddeyalli kanavarisidanu.} — Last night my younger brother talked in his sleep.
  \item \emph{Recall:} say the Kannada for to beg, then the Kannada for to cover oneself, then say \emph{kanavarisu} again
\end{itemize}

\begin{tcolorbox}[breakable,colback=teal!4,colframe=teal!35!black,title={Before you move on}]
What does \kn{ಕನವರಿಸು} mean? (\textquotedblleft{}To talk in sleep\textquotedblright{}.) Say it once more.
\end{tcolorbox}
\section[\texorpdfstring{\kn{ಹೊರಳು} (horaḷu)}{ಹೊರಳು (horaḷu)}]{\kn{ಹೊರಳು} (horaḷu) — to roll over; to toss}

\label{lesson:KA-C506-horalu}



\begin{quote}
Before the new one: say the Kannada for to cover oneself, then the Kannada for to talk in sleep.
\end{quote}

\subsection*{You'll want to know: \kn{ಹೊರಳು}}
\textbf{\kn{ಹೊರಳು}} — \emph{horaḷu} — \textquotedblleft{}to roll over; to toss\textquotedblright{}.

\begin{quote}
\kn{ಮಗು} \kn{ನಿದ್ದೆಯಲ್ಲಿ} \kn{ಹೊರಳಿತು}.
\end{quote}

\begin{quote}
\emph{magu niddeyalli horaḷitu.}
\end{quote}

\begin{quote}
The baby rolled over in its sleep.
\end{quote}

Past (I …): \kn{ಹೊರಳಿದೆ}, \emph{horaḷide}.

\kn{ಹೊರಳಾಡು}, to toss and turn all night.

\begin{cousinweb}
\textbf{\kn{ಹೊರಳು}} (\emph{horaḷu}) is native Dravidian: DEDR 4285 lists Kannada \emph{poraḷ}, \textquotedblleft{}to roll, wallow\textquotedblright{}. Old p- became h-.
\end{cousinweb}

\begin{grammarlens}[title={What you've built}]
One more word for describing what you lived and felt.
\end{grammarlens}

\subsection*{Your turn}
\begin{itemize}
\raggedright
  \item \emph{Say it:} \emph{horaḷu}
  \item \emph{Say it:} \emph{horaḷu}, once more
  \item \emph{Say it:} \emph{magu niddeyalli horaḷitu.} — The baby rolled over in its sleep.
  \item \emph{Recall:} say the Kannada for to cover oneself, then the Kannada for to talk in sleep, then say \emph{horaḷu} again
\end{itemize}

\begin{tcolorbox}[breakable,colback=teal!4,colframe=teal!35!black,title={Before you move on}]
What does \kn{ಹೊರಳು} mean? (\textquotedblleft{}To roll over; to toss\textquotedblright{}.) Say it once more.
\end{tcolorbox}
\section[\texorpdfstring{\kn{ಮೈಮುರಿ} (maimuri)}{ಮೈಮುರಿ (maimuri)}]{\kn{ಮೈಮುರಿ} (maimuri) — to stretch}

\label{lesson:KA-C506-maimuri}



\begin{quote}
Before the new one: say the Kannada for to talk in sleep, then the Kannada for to roll over.
\end{quote}

\subsection*{You'll want to know: \kn{ಮೈಮುರಿ}}
\textbf{\kn{ಮೈಮುರಿ}} — \emph{maimuri} — \textquotedblleft{}to stretch\textquotedblright{}.

\begin{quote}
\kn{ಬೆಳಗ್ಗೆ} \kn{ಎದ್ದ} \kn{ಬೆಕ್ಕು} \kn{ಮೈಮುರಿಯಿತು}.
\end{quote}

\begin{quote}
\emph{beḷagge edda bekku maimuriyitu.}
\end{quote}

\begin{quote}
The cat got up in the morning and stretched.
\end{quote}

Past (I …): \kn{ಮೈಮುರಿದೆ}, \emph{maimuride}.

\kn{ಮೈ}, the body + \kn{ಮುರಿ}, to bend: the long stretch after sleep.

\begin{cousinweb}
\textbf{\kn{ಮೈಮುರಿ}} (\emph{maimuri}) is a native Dravidian compound: \emph{mai}, \textquotedblleft{}the body\textquotedblright{} (DEDR 5073), and \emph{muri}, \textquotedblleft{}to bend; to stretch oneself with windings of the limbs\textquotedblright{} (DEDR 4977).
\end{cousinweb}

\begin{grammarlens}[title={What you've built}]
One more word for describing what you lived and felt.
\end{grammarlens}

\subsection*{Your turn}
\begin{itemize}
\raggedright
  \item \emph{Say it:} \emph{maimuri}
  \item \emph{Say it:} \emph{maimuri}, once more
  \item \emph{Say it:} \emph{beḷagge edda bekku maimuriyitu.} — The cat got up in the morning and stretched.
  \item \emph{Recall:} say the Kannada for to talk in sleep, then the Kannada for to roll over, then say \emph{maimuri} again
\end{itemize}

\begin{tcolorbox}[breakable,colback=teal!4,colframe=teal!35!black,title={Before you move on}]
What does \kn{ಮೈಮುರಿ} mean? (\textquotedblleft{}To stretch\textquotedblright{}.) Say it once more.
\end{tcolorbox}
\section[\texorpdfstring{\kn{ರೆಪ್ಪೆ} (reppe)}{ರೆಪ್ಪೆ (reppe)}]{\kn{ರೆಪ್ಪೆ} (reppe) — an eyelid; an eyelash}

\label{lesson:KA-C506-reppe}



\begin{quote}
Before the new one: say the Kannada for to roll over, then the Kannada for to stretch.
\end{quote}

\subsection*{You'll want to know: \kn{ರೆಪ್ಪೆ}}
\textbf{\kn{ರೆಪ್ಪೆ}} — \emph{reppe} — \textquotedblleft{}an eyelid; an eyelash\textquotedblright{}.

\begin{quote}
\kn{ಮಗುವಿನ} \kn{ರೆಪ್ಪೆಗಳು} \kn{ನಿಧಾನವಾಗಿ} \kn{ಮುಚ್ಚಿದವು}.
\end{quote}

\begin{quote}
\emph{maguvina reppegaḷu nidhānavāgi muccidavu.}
\end{quote}

\begin{quote}
The baby's eyelids slowly closed.
\end{quote}

Plural: \kn{ರೆಪ್ಪೆಗಳು}, \emph{reppegaḷu}.

\kn{ಕಣ್ಣುರೆಪ್ಪೆ} is the full form.

\begin{cousinweb}
\textbf{\kn{ರೆಪ್ಪೆ}} (\emph{reppe}) is native Dravidian: DEDR 5169 lists Kannada \emph{ṟeppe}, \textquotedblleft{}eyelid\textquotedblright{}, beside Tamil \emph{imai} forms and Telugu \emph{ṟeppa}.
\end{cousinweb}

\begin{grammarlens}[title={What you've built}]
One more word for describing what you lived and felt.
\end{grammarlens}

\subsection*{Your turn}
\begin{itemize}
\raggedright
  \item \emph{Say it:} \emph{reppe}
  \item \emph{Say it:} \emph{reppe}, once more
  \item \emph{Say it:} \emph{maguvina reppegaḷu nidhānavāgi muccidavu.} — The baby's eyelids slowly closed.
  \item \emph{Recall:} say the Kannada for to roll over, then the Kannada for to stretch, then say \emph{reppe} again
\end{itemize}

\begin{tcolorbox}[breakable,colback=teal!4,colframe=teal!35!black,title={Before you move on}]
What does \kn{ರೆಪ್ಪೆ} mean? (\textquotedblleft{}An eyelid; an eyelash\textquotedblright{}.) Say it once more.
\end{tcolorbox}
